\section{结论}

本文在通用纳米颗粒基准比较、纳米硒专用建模、增强特征筛选和多模态融合的早期研究基础上，加入了文献分组验证与独立 HUVEC 实验评价。来自 43 篇文献的 383 条早期记录支持了模型与特征组合的探索；另一份归档工作簿中的 434 条记录用于严格分组评价，两者的数据纳入、预处理和指标口径不同，不能直接合并或将其性能差异全部归因于某一验证策略。

早期含 IC50 的增强 XGBoost 取得三折交叉验证 MAE 7.20 个百分点和全量拟合 $R^2=0.957$，该结果保留为响应相关辅助分析。后续不含实测 IC50 的随机森林基础模型在文献分组嵌套评价中的折外 MAE 为 19.11 个百分点，当前化学增强、剂量—时间交互及固定预算迁移未显示稳定改善。因此，材料化学描述符具有继续采集与检验的价值，但现有证据不足以证明其已带来可迁移的独立收益，也不足以根据历史拟合指标确定部署模型。

新增九种 SeNPs 的 351 个 HUVEC 条件报告值，将原有外部验证计划推进到逐条件预测—实测比较。冻结模型的总体 MAE 为 14.51 个百分点、RMSE 为 20.14 个百分点、$R^2=0.446$，但对较长暴露时间的误差增大，预测材料排序与实测一致性较低。实验的 27 条曲线中 19 条可估计表观 IC50、8 条未达到阈值，而预测曲线均未跨越 50\% 存活率。SHAP 分析说明模型如何使用粒径和 Zeta 相关特征，不能替代材料性质的因果验证。

研究由此形成了从文献整理、模型比较和特征工程到严格分组评价及实验检验的连续路径。下一步需要核对浓度计量与换算基准、补齐实验流程和重复数据，并在独立材料或批次中检验模型冻结、预测与解盲过程。在这些条件得到验证前，当前模型适合作为进一步研究的基线，尚不能据此直接判定新材料安全性、可靠毒性排序或有限 IC50。
